%!TEX TS-program = xelatex
\documentclass[11pt, a4paper]{article}

\usepackage[fontset=windows]{ctex}
\usepackage{amsmath, amssymb, amsthm}
\usepackage{graphicx}
\usepackage{hyperref}
\usepackage{bookmark}
\usepackage{enumitem}

\newtheorem{problem}{题}
\renewcommand{\proofname}{\textbf{解答}}
\usepackage[margin=1in]{geometry}

\title{算法设计与分析 作业六}
\author{xxxxxxxxx xxx}
\date{\today}

\begin{document}

\maketitle

\begin{problem}[8.1]
\end{problem}

\begin{proof}
    设 $G_1$ 和 $G_2$ 的邻接矩阵分别为 $A=(a_{ij})$ 和 $B=(b_{ij})$。由于两个图的顶点集
    相同，并且输入已经按同一顶点顺序给出，只需检查任意不同顶点对在两个图中是否恰好有
    一条边。

    算法如下：对所有 $1\le i<j\le n$，比较 $a_{ij}$ 和 $b_{ij}$。若存在
    $a_{ij}=b_{ij}$，则输出 ``No''；若所有这样的顶点对都满足 $a_{ij}\ne b_{ij}$，则输出
    ``Yes''。显然算法只做
    \[
        \frac{n(n-1)}{2}
    \]
    次比较，故最坏情况下时间复杂度为 $\Theta(n^2)$。

    下面说明该阶也是下界。任意判定算法若没有检查某个顶点对 $(i,j)$，则对两个只在该位置
    不同的输入无法区分：一种输入中两个矩阵在该位置互补，另一种输入中两个矩阵在该位置
    相同，而算法的执行过程完全一样，却应给出不同答案。因此任一算法在最坏情况下至少要
    检查所有 $\frac{n(n-1)}{2}$ 个顶点对，时间复杂度下界为 $\Omega(n^2)$。
\end{proof}

\vspace{1em}

\begin{problem}[8.2]
\end{problem}

\begin{proof}
    采用 Horner 法从高次项到低次项依次计算。令
    \[
        P_1(x)=a_n,
    \]
    并对 $i=2,3,\ldots,n+1$ 定义
    \[
        P_i(x)=a_{n-i+1}+xP_{i-1}(x).
    \]
    最后得到
    \[
        P_{n+1}(x)=a_0+a_1x+\cdots+a_nx^n=P(x).
    \]
    每一步只需要一次乘法和一次加法，所以总时间为 $\Theta(n)$，占用额外空间为 $O(1)$。

    证明最优性。考虑任一算法，若某个系数 $a_i$ 在计算过程中完全没有被处理，则把输入中
    的 $a_i$ 改成另一个值后，算法执行过程不变，但由于题设 $x\ne0$，多项式值一般会发生
    改变，算法不可能同时正确。因此每个系数都至少需要被访问一次，时间下界为 $\Omega(n)$。
    上述算法达到该下界，故在最坏情况下是最优的。
\end{proof}

\vspace{1em}

\begin{problem}[8.4]
\end{problem}

\begin{proof}
    \begin{enumerate}[label=(\arabic*)]
        \item 若顺序归并，先令 $L=L_1$，再依次执行
        \[
            L\leftarrow L\cup L_j,\qquad j=2,3,\ldots,k,
        \]
        其中每次用通常的二路归并。第 $j$ 次归并时，前面已经合并出的表长为
        $(j-1)n/k$，新表长为 $n/k$，所以比较次数为 $O(jn/k)$。于是总时间为
        \[
            \sum_{j=2}^k O\left(\frac{jn}{k}\right)=O(nk).
        \]
        \item 更好的做法是二分归并。每轮把当前的有序表两两归并，表的个数约减半，直到
        只剩一个表。每一轮所有元素总共只参加一次归并，所以一轮时间为 $O(n)$；轮数为
        $\lceil\log k\rceil$，故最坏情况下时间复杂度为 $O(n\log k)$。
        \item 证明下界。把所有输入按比较过程形成决策树。每个输入表有 $n/k$ 个元素，保持
        每个表内部顺序不变时，可能的总排列数为
        \[
            \frac{n!}{((n/k)!)^k}.
        \]
        因而决策树至少有这么多片叶子，深度至少为
        \[
            \log\frac{n!}{((n/k)!)^k}
            =\Theta(n\log n-k\cdot (n/k)\log(n/k))
            =\Theta(n\log k).
        \]
        因此任何基于比较的算法最坏情况下都需要 $\Omega(n\log k)$ 次比较。二分归并达到
        该下界，所以是渐近最优的。
    \end{enumerate}
\end{proof}

\vspace{1em}

\begin{problem}[8.9]
\end{problem}

\begin{proof}
    \begin{enumerate}[label=(\arabic*)]
        \item 先用线性时间选择算法找出 $L$ 中第 $m$ 小的元素 $a$。再扫描一遍 $L$，把所有
        小于 $a$ 的元素收集出来。由于 $m=o(n/\log n)$，对这 $m-1$ 个元素排序需要
        \[
            O(m\log m)=o(n)
        \]
        时间。最后按升序输出这些元素以及 $a$。总时间为
        \[
            O(n)+o(n)=O(n).
        \]

        \item 若存在一个 $O(n)$ 时间算法能在 $m=\omega(n/\log n)$ 时完成该任务，则它也能
        在 $O(n)$ 时间内把任意给定的 $m$ 个数排序：把这 $m$ 个数与另外 $n-m$ 个都比它们
        大的辅助数放在同一表中，运行该算法，输出的前 $m$ 小元素正好就是原来 $m$ 个数的
        有序序列。可是基于比较的排序需要 $\Omega(m\log m)$ 时间，而当
        $m=\omega(n/\log n)$ 时有
        \[
            m\log m=\omega(n),
        \]
        与 $O(n)$ 矛盾。因此此时不存在基于比较的 $O(n)$ 时间算法。
    \end{enumerate}
\end{proof}

\vspace{1em}

\begin{problem}[8.10]
\end{problem}

\begin{proof}
    设要求 $n$ 个互异数中的第 $k$ 小元素。为了确定答案，算法最终必须知道至少
    $k-1$ 个元素小于它，并且至少 $n-k$ 个元素大于它。先考虑小于答案的元素。若某个小于
    答案的元素从未在比较中输给任何元素，则它仍可能是第 $k$ 小元素，答案无法确定；因此
    至少要通过 $k-1$ 次比较给这些元素各分配一次 ``小于答案'' 的证据。同理，为了排除大于
    答案的元素，至少需要 $n-k$ 次比较给它们各分配一次 ``大于答案'' 的证据。

    另外，答案元素本身必须至少参加
    \[
        \min\{k-1,n-k\}
    \]
    次比较。否则，与它比较过的较小元素不足 $k-1$ 个，或较大元素不足 $n-k$ 个，总可以
    构造另一组与已有比较结果一致的输入，使答案发生改变。

    因此任何基于比较的算法最坏情况下至少需要
    \[
        (k-1)+(n-k)+\min\{k-1,n-k\}
        =n+\min\{k,n-k+1\}-2
    \]
    次比较。这就是所要求的下界。
\end{proof}

\end{document}
